\documentclass{article}

\usepackage[utf8]{inputenc}
\usepackage{helvet}
\usepackage{geometry}
\usepackage{xcolor}
\usepackage{eso-pic}
\usepackage{fancyhdr}
\usepackage{parskip}
\usepackage{microtype}
\usepackage[english, ukrainian]{babel}
\usepackage{url}
\usepackage{amsmath}
\usepackage{amssymb}
\usepackage{amsthm}
\usepackage{longtable}
\usepackage{booktabs}
\usepackage{hyperref}
\usepackage{titlesec}

\definecolor{nato}{HTML}{293757}
\definecolor{footergray}{gray}{0.65}

\geometry{a4paper,left=3cm,right=3cm,top=3cm,bottom=3cm}
\renewcommand{\familydefault}{\sfdefault}
\setcounter{secnumdepth}{3}

\titleformat{\section}
  {\color{nato}\Huge\bfseries}{\thesection.}{1em}{}
\titlespacing*{\section}{0pt}{30pt}{12pt}

\titleformat{\subsection}
  {\color{nato}\LARGE\bfseries}{\thesubsection.}{1em}{}
\titlespacing*{\subsection}{0pt}{20pt}{8pt}

\titleformat{\subsubsection}
  {\color{nato}\Large\bfseries}{\thesubsubsection.}{1em}{}
\titlespacing*{\subsubsection}{0pt}{10pt}{6pt}

\fancyhf{}
\fancyhead[R]{\small\color{footergray} 27 червня 2026}
\fancyfoot[L]{\small\color{footergray} Техноробочий проект ERP/1. Модуль Освіта}
\fancyfoot[R]{\small\color{footergray}\thepage}

\newcommand\HUGE[1]{\fontsize{40}{40}\selectfont #1}

\renewcommand{\headrulewidth}{0pt}
\renewcommand{\footrulewidth}{0.4pt}
\renewcommand{\footrule}{\color{footergray}\hrule width \textwidth height 0.4pt}

\begin{document}

\pagestyle{fancy}

\begin{titlepage}
  \AddToShipoutPictureBG*{\AtPageLowerLeft{\color{nato}\rule{\paperwidth}{\paperheight}}}
  \color{white}\thispagestyle{empty}\vspace*{3cm}\centering
  {\Huge  \textbf{ERP/1: Освіта} \par} \vspace{1.5cm}
  {\HUGE  \textbf{Техноробочий проект} \par} \vspace{1.5cm}
  {\LARGE \textbf{на створення та впровадження модуля «Освіта» \\ як засобу інформатизації} \par} \vspace{1.5cm}
  {\large Згідно з ДСТУ 34.201-89 та Постановою КМУ № 205 від 21.02.2025 \par}
  {\large Формалізація вимог за стандартом ISO/IEC/IEEE 29148:2018 \par}
  \vfill
  {\large 2026 \copyright\ Системи Електронного Урядування \par}
\end{titlepage}

\newpage

\begin{center}
  {\Huge\color{nato}\bfseries Зміст\par}
\end{center}
\vspace{40pt}
\tableofcontents

\newpage

\section{Загальні відомості про засіб інформатизації}

\subsection{1.1. Найменування та підстави для розробки}
\textbf{Найменування засобу інформатизації:} Модуль «Освіта» інформаційної системи управління підприємством ERP/1 (далі --- Модуль «Освіта»).
\\
\textbf{Відомості про замовника:} ТОВ "Системи Електронного Урядування" в ТОВ "Криптографічні телесистеми"..
\\
\textbf{Відомості про виконавця робіт:} Визначається замовником на конкурсній основі.
\\
\textbf{Підстави для виконання робіт:}
\begin{itemize}
    \item Закон України «Про освіту» (із змінами та доповненнями);
    \item Закон України «Про захист персональних даних»;
    \item Постанова Кабінету Міністрів України від 21 лютого 2025 року № 205 «Деякі питання створення, адміністрування та забезпечення функціонування засобу інформатизації»;
    \item Програма інформатизації ERP/1.
\end{itemize}

\subsection{1.2. Призначення та цілі створення засобу}
Модуль «Освіта» призначений для автоматизації та управління освітніми процесами у закладах освіти різного рівня (ЗСО, ПО, ФПО, ВО).
\\
\textbf{Цілі створення засобу:}
\begin{itemize}
    \item Створення сучасного середовища ведення освітньої документації в електронному вигляді з накладанням КЕП;
    \item Автоматизація розробки та затвердження навчальних планів і робочих програм;
    \item Ведення електронного журналу успішності та відвідуваності в режимі реального часу;
    \item Контроль виконання бізнес-правил освітнього процесу та валідація оцінок;
    \item Автоматизоване планування розкладу занять на основі ресурсних обмежень;
    \item Забезпечення інтеграції з Єдиною державною електронною базою з питань освіти (ЄДЕБО).
\end{itemize}

\subsection{1.3. План-графік виконання етапів робіт}
Роботи з розробки та впровадження Модуля «Освіта» розбиваються на такі етапи:
\begin{enumerate}
    \item \textbf{Етап 1: Технічне проектування} --- погодження UX/UI прототипів у Figma, уточнення структури даних Mnesia. (Тривалість: 1 місяць).
    \item \textbf{Етап 2: Розробка інформаційного забезпечення та схем} --- опис та реалізація Erlang-рекордів, створення базових таблиць Mnesia. (Тривалість: 1.5 місяці).
    \item \textbf{Етап 3: Реалізація бізнес-процесів} --- написання BPE DSL та інтеграція з рушієм процесів. (Тривалість: 2 місяці).
    \item \textbf{Етап 4: Інтеграція з КЕП та ЄДЕБО} --- реалізація криптографічних підписів та API обміну. (Тривалість: 1 місяць).
    \item \textbf{Етап 5: Комплексне тестування} --- автотестування (Unit, Integration), UAT приймання. (Тривалість: 1 місяць).
    \item \textbf{Етап 6: Впровадження та дослідна експлуатація}. (Тривалість: 1.5 місяці).
\end{enumerate}

\newpage
\section{Відомості про робочий процес та умови експлуатації}

\subsection{2.1. Опис бізнес-процесів (BPMN / FSM)}
Освітні бізнес-процеси в Модулі «Освіта» реалізуються за допомогою скінченних автоматів (FSM) та інтегруються з рушієм процесів BPE:

\subsubsection{2.1.1. Процес оцінювання (Student Assessment FSM)}
Процес охоплює виставлення поточних та підсумкових балів, розрахунок компетенцій та затвердження КЕП:
\begin{enumerate}
    \item \textbf{StartAssessment:} Початок періоду оцінювання (початок уроку або семестрової атестації).
    \item \textbf{EnterGrades:} Педагог вводить оцінки у систему.
    \item \textbf{ValidateGrades:} Автоматична перевірка введених оцінок за допомогою валідатора \texttt{grade\_}\-\texttt{validator.erl} на відповідність порогам. При успішній перевірці перехід до \texttt{CalculateCompetences}, інакше --- повернення до \texttt{EnterGrades}.
    \item \textbf{CalculateCompetences:} Автоматичний розрахунок середнього балу та інтегральних показників компетенцій за темою/семестром.
    \item \textbf{TeacherSign:} Педагог підписує журнал КЕП. Якщо оцінки містять незадовільні балі ($<4$) при підсумковій атестації, процес вимагає переходу до стану \texttt{ViceCheck}.
    \item \textbf{ViceCheck:} Завуч/методист перевіряє та погоджує відомість.
    \item \textbf{AssessmentSigned:} Завершення процесу, фіксація записів.
\end{enumerate}

\subsubsection{2.1.2. Процес проведення занять (Lesson Execution)}
\begin{enumerate}
    \item \textbf{LessonStart:} Початок заняття за затвердженим розкладом.
    \item \textbf{TakeAttendance:} Фіксація відсутності учнів (Present/Absent).
    \item \textbf{AssessStudents:} Виставлення оцінок за роботу на уроці та домашні завдання.
    \item \textbf{CalculateAverages:} Розрахунок поточної середньої успішності.
    \item \textbf{SignJournal:} Накладання КЕП на журнал протягом 48-годинного вікна.
    \item \textbf{LessonClosed:} Фіксація журналу, архівування.
\end{enumerate}

\subsubsection{2.1.3. Затвердження навчального плану (Curriculum Approval)}
\begin{enumerate}
    \item \textbf{Start:} Створення проекту навчального плану методистом.
    \item \textbf{CreatePlan:} Наповнення плану дисциплінами, годинами та кредитами.
    \item \textbf{MethodistReview:} Перевірка плану на відповідність лімітам освітньої програми.
    \item \textbf{ViceApprove:} Погодження плану проректором/завучем.
    \item \textbf{DirectorApprove:} Остаточне затвердження КЕП директора закладу.
    \item \textbf{GenerateSchedule:} Перехід до автоматичного планування розкладу занять.
    \item \textbf{PlanApproved:} Активація плану.
\end{enumerate}

\subsubsection{2.1.4. Генерація розкладу занять (Schedule Generation)}
\begin{enumerate}
    \item \textbf{PlanApproved:} Запуск генерації після затвердження плану.
    \item \textbf{LoadPlanAndConstraints:} Завантаження плану та ресурсних обмежень (кількість аудиторій, доступність викладачів).
    \item \textbf{GenerateSchedule:} Робота евристичного алгоритму генерації. При конфліктах --- перехід до \texttt{ManualAdjustment}, при успіху --- до \texttt{ViceDirectorApprove}.
    \item \textbf{ManualAdjustment:} Ручне редагування розкладу диспетчером.
    \item \textbf{ViceDirectorApprove:} Затвердження розкладу керівництвом закладу.
    \item \textbf{PublishAndNotify:} Публікація та розсилка повідомлень про розклад.
\end{enumerate}

\subsection{2.2. Опис ролей та прав доступу (ABAC)}
Модель контролю доступу базується на атрибутах (ABAC) та перевіряється на рівні ядра авторизації.

\paragraph{Матриця базового доступу:}
\begin{itemize}
    \item \textbf{Педагог (Teacher):} Доступ \texttt{read}, \texttt{write}, \texttt{sign} для журналів (\texttt{JournalRecord}), якщо він є викладачем цього курсу або курс входить до його списку дисциплін. Редагування обмежено 48 годинами з моменту заняття.
    \item \textbf{Учень (Student):} Доступ \texttt{read} тільки до власних оцінок. Заборонено \texttt{write}, \texttt{sign}, \texttt{delete}.
    \item \textbf{Батько/Мати (Parent):} Доступ \texttt{read} для оцінок дитини (згідно зі списком зв'язків).
    \item \textbf{Адміністратор/Завуч (Admin/Vice):} Перегляд усіх документів, затвердження відомостей, редагування планів та довідників.
\end{itemize}

\paragraph{Формальні логічні правила ABAC:}
\begin{itemize}
    \item \textbf{Правило 1 (Педагог):}
    $$\text{subject.roles contains 'Teacher'} \land \text{object.entity\_type == 'JournalRecord'}$$
    $$\land\ (\text{object.teacher\_id == subject.person\_id} \lor \text{object.course\_id in subject.courses\_taught})$$
    $$\rightarrow \text{Permit: read, write, sign.}$$
    \item \textbf{Правило 2 (Учень):}
    $$\text{subject.roles contains 'Student'} \land \text{object.entity\_type == 'JournalRecord'} \land \text{object.student\_id == subject.person\_id}$$
    $$\rightarrow \text{Permit: read (тільки власні оцінки). Deny: write, sign.}$$
    \item \textbf{Правило 3 (Батько):}
    $$\text{subject.roles contains 'Parent'} \land \text{object.student\_id in subject.children\_ids}$$
    $$\rightarrow \text{Permit: read. Deny: write, sign.}$$
    \item \textbf{Правило 4 (КЕП):}
    $$\text{subject.kep\_cert\_thumbprint is valid} \land \text{subject.roles contains 'Teacher'} \land \text{object.status == 'draft'}$$
    $$\rightarrow \text{Permit: sign.}$$
    \item \textbf{Правило 5 (Час редагування):}
    $$\text{subject.roles contains 'Teacher'} \land \text{env.time} \le \text{lesson.start\_time} + 48\text{ годин} \land \text{object.status == 'draft'}$$
    $$\rightarrow \text{Permit: write.}$$
\end{itemize}

\subsection{2.3. Умови експлуатації та системне середовище}
Модуль «Освіта» розробляється як хмарна веб-платформа.
\\
\textbf{Системне середовище:}
\begin{itemize}
    \item \textbf{Середовище виконання:} Erlang/OTP версії 26 або вище.
    \item \textbf{База даних:} Вбудована СУБД Erlang Mnesia. Дозволяється використання будь-яких сторонніх SQL СУБД (PostgreSQL, MySQL, SQLite) для підвищення продуктивності та відмовостійкості, реплікації та аналітики.
    \item \textbf{Веб-сервер:} Cowboy (вбудований в Erlang-додаток).
    \item \textbf{Фронтенд-платформа:} Веб-браузери з підтримкою HTML5, CSS3, JavaScript та Secure WebSockets.
\end{itemize}

\newpage
\section{Інформаційне та програмне забезпечення}

\subsection{3.1. Інформаційне забезпечення (Схеми та Дані)}

\subsubsection{3.1.1. Визначення структур даних (Erlang Records)}
Вся інформаційна модель бази даних описується у вигляді рекорді Erlang:
\begin{verbatim}
-record(curriculum, {
    id,                  % UUID (Primary Key)
    program_id,          % UUID (FK -> EducationalProgram)
    group_id,            % UUID (FK -> Group)
    academic_year,       % Integer (навчальний рік)
    semester,            % Integer (1-2)
    total_credits,       % Decimal (загальна кількість кредитів)
    total_hours,         % Integer (загальна кількість годин)
    approved_by,         % UUID (FK -> Teacher, КЕП підпис)
    approved_at,         % Timestamp
    status = draft,      % draft | pending | approved | archive
    form_of_study_id,    % UUID (FK -> FormOfStudy)
    language,            % Enum (ua, en)
    is_active = true,    % Boolean
    metadata = #{}       % Map для додаткових неструктурованих даних
}).

-record(curriculum_item, {
    id,                  % UUID (PK)
    curriculum_id,       % UUID (FK -> Curriculum)
    course_id,           % UUID (FK -> Course)
    semester,            % Integer
    hours_lectures,      % Integer
    hours_practice,      % Integer
    hours_lab,           % Integer
    hours_self,          % Integer
    credits_ects,        % Decimal
    grade_type_id,       % UUID (FK -> GradeType)
    teacher_id,          % UUID (FK -> Teacher)
    order,               % Integer
    is_mandatory = true  % Boolean
}).

-record(form_of_study, {
    id,                  % UUID (PK)
    code,                % String (FULL, DIST, MIX)
    name_ua,             % String (українська назва)
    name_en,             % String (англійська назва)
    description,         % Text
    is_active = true,    % Boolean
    order,               % Integer
    supports_mixed = false,
    requires_online_platform = false,
    legal_basis,
    metadata = #{}
}).

-record(grade_type, {
    id,                  % UUID (PK)
    code,                % String
    name_ua,             % String
    name_en,             % String
    scale_type,          % numeric_12 | ects | competence | score_100
    min_value,           % Decimal
    max_value,           % Decimal
    passing_threshold,   % Decimal
    weight = 0.0,        % Decimal (0.0 - 1.0)
    is_final = false,    % Boolean
    requires_approval = false,
    metadata = #{}
}).

-record(educational_program_type, {
    id,
    code,
    name_ua,
    name_en,
    education_level,
    program_category,
    duration_years_min,
    duration_years_max,
    ects_min,
    description,
    is_active = true,
    metadata = #{}
}).

-record(educational_program, {
    id, code, name_ua, name_en, level, specialty_code, specialty_name,
    qualification, duration_years, credits_ects, form_of_study_id,
    language, accreditation_status, accreditation_until,
    license_number, department_id, is_active = true, metadata = #{}
}).
\end{verbatim}

\subsubsection{3.1.2. Опис довідників та реєстрових сутностей}

\paragraph{Словник: Типи освітніх програм (EducationalProgramType)}
\begin{longtable}{p{3.5cm}p{3cm}p{2cm}p{5.5cm}}
\toprule
\textbf{Реквізит} & \textbf{Тип даних} & \textbf{Обов'язковий} & \textbf{Опис} \\
\midrule
id & UUID & Так & Унікальний ідентифікатор \\
code & String (30) & Так & Код типу програми (GENERAL\_SECONDARY тощо) \\
name\_ua & String (300) & Так & Назва українською \\
duration\_years\_min & Decimal & Так & Мінімальна тривалість (років) \\
ects\_min & Integer & Так & Мінімальний обсяг кредитів ECTS \\
is\_active & Boolean & Так & Чи є активним запис \\
\bottomrule
\end{longtable}

\paragraph{Словник: Форми навчання (FormOfStudy)}
\begin{longtable}{p{3.5cm}p{3cm}p{2cm}p{5.5cm}}
\toprule
\textbf{Реквізит} & \textbf{Тип даних} & \textbf{Обов'язковий} & \textbf{Опис} \\
\midrule
id & UUID & Так & Унікальний ідентифікатор \\
code & String (10) & Так & Код форми навчання (FULL, DIST, MIX) \\
name\_ua & String (150) & Так & Назва українською \\
supports\_mixed & Boolean & Так & Підтримка змішаних форматів \\
requires\_online\_platform& Boolean & Так & Потреба в LMS платформі \\
\bottomrule
\end{longtable}

\paragraph{Реєстр: Особові справи (Student)}
\begin{longtable}{p{3.5cm}p{3cm}p{2cm}p{5.5cm}}
\toprule
\textbf{Реквізит} & \textbf{Тип даних} & \textbf{Обов'язковий} & \textbf{Опис} \\
\midrule
id & UUID & Так & Унікальний ідентифікатор \\
external\_id & String & Ні & Зв'язок з ЄДЕБО \\
full\_name & String & Так & Прізвище, Ім'я, По батькові \\
date\_of\_birth & Date & Так & Дата народження \\
group\_id & UUID & Так & Зв'язок з групою (Group.id) \\
status & Enum & Так & Поточний статус (Active, Graduated тощо) \\
kep\_cert & Binary & Ні & Збережений КЕП сертифікат \\
\bottomrule
\end{longtable}

\paragraph{Реєстр: Записи оцінювання (JournalRecord)}
\begin{longtable}{p{3.5cm}p{3cm}p{2cm}p{5.5cm}}
\toprule
\textbf{Реквізит} & \textbf{Тип даних} & \textbf{Обов'язковий} & \textbf{Опис} \\
\midrule
id & UUID & Так & Унікальний ідентифікатор \\
student\_id & UUID & Так & FK $\rightarrow$ Student \\
course\_id & UUID & Так & FK $\rightarrow$ Course \\
lesson\_date & Date & Так & Дата уроку \\
attendance\_status & Enum & Так & Присутність (Present, Absent, Sick) \\
grade & Decimal & Ні & Числова оцінка (1.0 - 12.0 або 0.0 - 100.0) \\
grade\_type & Enum & Так & Посилання на GradeType.code \\
competence\_scores & Map & Ні & Бали за компетенції \\
signed\_by & UUID & Ні & Вчитель, що підписав відомість КЕП \\
signed\_at & Timestamp & Ні & Дата підпису \\
\bottomrule
\end{longtable}

\subsubsection{3.1.3. Алгоритми конвертації оцінок}
Модуль виконує автоматичний перерахунок оцінок за такими правилами:
\begin{itemize}
    \item \textbf{12-бальна шкала у літерну ECTS:}
    $$\text{ECTS} = \begin{cases} 
      A, & \text{якщо } \text{Grade} \ge 10 \\
      B, & \text{якщо } \text{Grade} == 9 \\
      C, & \text{якщо } \text{Grade} == 8 \\
      D, & \text{якщо } \text{Grade} == 7 \\
      E, & \text{якщо } \text{Grade} == 6 \\
      FX, & \text{якщо } \text{Grade} \ge 4 \land \text{Grade} < 6 \\
      F, & \text{якщо } \text{Grade} < 4 
   \end{cases}$$
    \item \textbf{12-бальна шкала у відсоткову 100-бальну:}
    $$\text{Percent} = \text{round}\left( \frac{\text{Grade}}{12} \times 100 \right)$$
    \item \textbf{Конвертація компетенцій у числову шкалу:}
    Розвинена $\rightarrow 12$; Базова $\rightarrow 8$; Початкова $\rightarrow 5$; Не сформована $\rightarrow 2$.
\end{itemize}

\subsection{3.2. Програмне забезпечення (Реалізація)}

\subsubsection{3.2.1. bpe\_education\_assessment.erl (DSL процесу оцінювання)}
Код модуля опису FSM процесу оцінювання та валідації успішності студентів:
\begin{verbatim}
-module(bpe_education_assessment).
-include("bpe.hrl").
-compile(export_all).

def() ->
    #process{
        name = 'Student Assessment',
        beginEvent = 'StartAssessment',
        endEvent = 'AssessmentSigned',
        tasks = [
            #beginEvent{name='StartAssessment'},
            #userTask{name='EnterGrades', module=?MODULE},
            #serviceTask{name='ValidateGrades', module=?MODULE},
            #serviceTask{name='CalculateCompetences', module=?MODULE},
            #userTask{name='TeacherSign', module=?MODULE},
            #userTask{name='ViceCheck', module=?MODULE},
            #endEvent{name='AssessmentSigned'}
        ],
        flows = [
            #sequenceFlow{name='1', source='StartAssessment', target='EnterGrades'},
            #sequenceFlow{name='2', source='EnterGrades', target='ValidateGrades'},
            #sequenceFlow{name='3', source='ValidateGrades', 
                         target='CalculateCompetences', condition="validation_ok"},
            #sequenceFlow{name='4', source='ValidateGrades', 
                         target='EnterGrades', condition="validation_failed"},
            #sequenceFlow{name='5', source='CalculateCompetences', target='TeacherSign'},
            #sequenceFlow{name='6', source='TeacherSign', 
                         target='ViceCheck', condition="requires_approval"},
            #sequenceFlow{name='7', source='TeacherSign', 
                         target='AssessmentSigned', condition="no_approval_needed"},
            #sequenceFlow{name='8', source='ViceCheck', 
                         target='AssessmentSigned', condition="vice_approved"}
        ],
        roles = [teacher, vice_rector]
    }.

action({serviceTask, 'ValidateGrades'}, Proc) ->
    case validate_grades(Proc) of
        ok -> {reply, Proc, "validation_ok"};
        _  -> {reply, Proc, "validation_failed"}
    end;

action({userTask, 'TeacherSign'}, Proc) ->
    case requires_vice_approval(Proc) of
        true -> {reply, Proc, "requires_approval"};
        false -> {reply, Proc, "no_approval_needed"}
    end;

action(_, Proc) -> {reply, Proc}.
\end{verbatim}

\subsubsection{3.2.2. bpe\_education\_lesson.erl (DSL процесу проведення заняття)}
\begin{verbatim}
-module(bpe_education_lesson).
-include("bpe.hrl").
-compile(export_all).

def() ->
    #process{
        name = 'Lesson Execution and Assessment',
        beginEvent = 'LessonStart',
        endEvent = 'LessonClosed',
        tasks = [
            #beginEvent{name='LessonStart'},
            #userTask{name='TakeAttendance', module=?MODULE},
            #userTask{name='AssessStudents', module=?MODULE},
            #serviceTask{name='CalculateAverages', module=?MODULE},
            #userTask{name='SignJournal', module=?MODULE},
            #endEvent{name='LessonClosed'}
        ],
        flows = [
            #sequenceFlow{name='1', source='LessonStart', target='TakeAttendance'},
            #sequenceFlow{name='2', source='TakeAttendance', target='AssessStudents'},
            #sequenceFlow{name='3', source='AssessStudents', target='CalculateAverages'},
            #sequenceFlow{name='4', source='CalculateAverages', 
                         target='SignJournal', condition="averages_ok"},
            #sequenceFlow{name='5', source='CalculateAverages', 
                         target='AssessStudents', condition="averages_need_correction"},
            #sequenceFlow{name='6', source='SignJournal', 
                         target='LessonClosed', condition="kep_signed"}
        ],
        events = [#timeoutEvent{name='JournalDeadline', 
                               timeout=#timeout{spec={hours,48}}}]
    }.

action({serviceTask, 'CalculateAverages'}, Proc) ->
    case averages_valid(Proc) of
        true -> {reply, Proc, "averages_ok"};
        false -> {reply, Proc, "averages_need_correction"}
    end;

action({userTask, 'SignJournal'}, Proc) ->
    case has_valid_kep(Proc) of
        true -> {reply, Proc, "kep_signed"};
        false -> {reply, Proc}
    end;

action(_, Proc) -> {reply, Proc}.
\end{verbatim}

\subsubsection{3.2.3. bpe\_curriculum\_approval.erl (DSL затвердження навчального плану)}
\begin{verbatim}
-module(bpe_curriculum_approval).
-include("bpe.hrl").
-compile(export_all).

def() ->
    #process{
        name = 'Curriculum Approval',
        beginEvent = 'Start',
        endEvent = 'PlanApproved',
        tasks = [
            #beginEvent{name='Start'},
            #userTask{name='CreatePlan', module=?MODULE},
            #userTask{name='MethodistReview', module=?MODULE},
            #userTask{name='ViceApprove', module=?MODULE},
            #userTask{name='DirectorApprove', module=?MODULE},
            #serviceTask{name='GenerateSchedule', module=?MODULE},
            #endEvent{name='PlanApproved'}
        ],
        flows = [
            #sequenceFlow{name='1', source='Start', target='CreatePlan'},
            #sequenceFlow{name='2', source='CreatePlan', target='MethodistReview'},
            #sequenceFlow{name='3', source='MethodistReview', 
                         target='ViceApprove', condition="methodist_ok"},
            #sequenceFlow{name='4', source='MethodistReview', 
                         target='CreatePlan', condition="methodist_reject"},
            #sequenceFlow{name='5', source='ViceApprove', 
                         target='DirectorApprove', condition="vice_ok"},
            #sequenceFlow{name='6', source='ViceApprove', 
                         target='MethodistReview', condition="vice_reject"},
            #sequenceFlow{name='7', source='DirectorApprove', 
                         target='GenerateSchedule', condition="director_ok"},
            #sequenceFlow{name='8', source='GenerateSchedule', target='PlanApproved'}
        ],
        roles = [methodist, vice_rector, director]
    }.

action({userTask, 'CreatePlan'}, Proc) -> {reply, Proc};
action({userTask, 'MethodistReview'}, Proc) ->
    case validate_curriculum(Proc) of
        ok -> {reply, Proc, "methodist_ok"};
        _  -> {reply, Proc, "methodist_reject"}
    end;

action({userTask, 'ViceApprove'}, Proc) ->
    case vice_approval(Proc) of
        ok -> {reply, Proc, "vice_ok"};
        _  -> {reply, Proc, "vice_reject"}
    end;

action({userTask, 'DirectorApprove'}, Proc) ->
    case director_kep_sign(Proc) of
        ok -> {reply, Proc, "director_ok"};
        _  -> {reply, Proc}
    end;

action({serviceTask, 'GenerateSchedule'}, Proc) ->
    {reply, Proc};

action(_, Proc) -> {reply, Proc}.
\end{verbatim}

\subsubsection{3.2.4. bpe\_schedule\_generation.erl (DSL генерації розкладу)}
\begin{verbatim}
-module(bpe_schedule_generation).
-include("bpe.hrl").
-compile(export_all).

def() ->
    #process{
        name = 'Schedule Generation',
        beginEvent = 'PlanApproved',
        endEvent = 'SchedulePublished',
        tasks = [
            #beginEvent{name='PlanApproved'},
            #serviceTask{name='LoadPlanAndConstraints', module=?MODULE},
            #serviceTask{name='GenerateSchedule', module=?MODULE},
            #userTask{name='ManualAdjustment', module=?MODULE},
            #userTask{name='ViceDirectorApprove', module=?MODULE},
            #serviceTask{name='PublishAndNotify', module=?MODULE},
            #endEvent{name='SchedulePublished'}
        ],
        flows = [
            #sequenceFlow{name='1', source='PlanApproved', target='LoadPlanAndConstraints'},
            #sequenceFlow{name='2', source='LoadPlanAndConstraints', target='GenerateSchedule'},
            #sequenceFlow{name='3', source='GenerateSchedule', 
                         target='ViceDirectorApprove', condition="no_conflicts"},
            #sequenceFlow{name='4', source='GenerateSchedule', 
                         target='ManualAdjustment', condition="has_conflicts"},
            #sequenceFlow{name='5', source='ManualAdjustment', target='GenerateSchedule'},
            #sequenceFlow{name='6', source='ViceDirectorApprove', 
                         target='PublishAndNotify', condition="approved"},
            #sequenceFlow{name='7', source='PublishAndNotify', target='SchedulePublished'}
        ]
    }.

action({serviceTask, 'LoadPlanAndConstraints'}, Proc) -> {reply, Proc};
action({serviceTask, 'GenerateSchedule'}, Proc) ->
    case generate_schedule_algorithm(Proc) of
        {ok, Schedule} -> 
            {reply, Proc#process{docs = [Schedule | Proc#process.docs]}, "no_conflicts"};
        {error, Conflicts} -> 
            {reply, Proc#process{docs = [Conflicts | Proc#process.docs]}, "has_conflicts"}
    end;
action({userTask, 'ManualAdjustment'}, Proc) -> {reply, Proc};
action({userTask, 'ViceDirectorApprove'}, Proc) ->
    case has_valid_kep(Proc) of
        true -> {reply, Proc, "approved"};
        false -> {reply, Proc}
    end;
action({serviceTask, 'PublishAndNotify'}, Proc) -> {reply, Proc};
action(_, Proc) -> {reply, Proc}.
\end{verbatim}

\subsubsection{3.2.5. grade\_validator.erl (Модуль валідації надійності даних та оцінок)}
Код Erlang-модуля для перевірки порогових значень та відсутності невалідності типів даних:
\begin{verbatim}
-module(grade_validator).
-export([validate/2, validate_thresholds/2]).

-record(grade_threshold, {
    id,
    grade_type_code,
    threshold_name,
    value,
    operator,
    message_ua,
    is_blocking = true,
    requires_approval = false
}).

validate(GradeTypeCode, Value) ->
    Thresholds = fetch_thresholds(GradeTypeCode),
    validate_thresholds(Value, Thresholds).

validate_thresholds(_Value, []) -> ok;
validate_thresholds(Value, [T|Rest]) ->
    case check_threshold(Value, T) of
        ok -> validate_thresholds(Value, Rest);
        {error, Msg} -> {error, Msg}
    end.

check_threshold(Value, #grade_threshold{operator = '>=', value = Th}) when Value >= Th -> ok;
check_threshold(Value, #grade_threshold{operator = '>', value = Th}) when Value > Th -> ok;
check_threshold(Value, #grade_threshold{operator = '<=', value = Th}) when Value =< Th -> ok;
check_threshold(_, T) -> {error, T#grade_threshold.message_ua}.

fetch_thresholds(<<"CURRENT_12">>) ->
    [
        #grade_threshold{
            id = 1,
            grade_type_code = <<"CURRENT_12">>,
            threshold_name = <<"passing">>,
            value = 6.0,
            operator = '>=',
            message_ua = <<"Poperedzhennia: Ocinka nyzhche prokhidnogo balu (6).">>,
            is_blocking = false
        },
        #grade_threshold{
            id = 2,
            grade_type_code = <<"CURRENT_12">>,
            threshold_name = <<"critical">>,
            value = 4.0,
            operator = '>=',
            message_ua = <<"Pomylka: Ocinka ie krytychno nezadovilnoiu. Potriben komentar pedagoga.">>,
            is_blocking = true
        }
    ];
fetch_thresholds(_) -> [].
\end{verbatim}

\newpage
\section{Вимоги до засобу за ISO/IEC/IEEE 29148}

\subsection{4.1. Функціональні вимоги}
\begin{itemize}
    \item \textbf{[REQ-EDU-FUN-001]} Модуль «Освіта» повинен забезпечувати ведення та оновлення словників типів освітніх програм (\texttt{EducationalProgramType}), форм навчання (\texttt{FormOfStudy}) та типів оцінок (\texttt{GradeType}) з можливістю версіонування записів.
    \item \textbf{[REQ-EDU-FUN-002]} Модуль «Освіта» повинен підтримувати автоматичний перерахунок оцінок між 12-бальною шкалою, літерними оцінками ECTS та 100-бальною відсотковою шкалою за заданими математичними формулами.
    \item \textbf{[REQ-EDU-FUN-003]} Модуль «Освіта» повинен реалізовувати скінченні автомати (FSM) життєвого циклу занять та ведення журналів.
    \item \textbf{[REQ-EDU-FUN-004]} Модуль «Освіта» повинен забезпечувати автоматичне блокування редагування записів електронного журналу через 48 годин після фактичного початку заняття.
    \item \textbf{[REQ-EDU-FUN-005]} Модуль «Освіта» повинен надавати викладачам інтерфейс для накладання кваліфікованого електронного підпису (КЕП) на відомості оцінювання та журнали.
    \item \textbf{[REQ-EDU-FUN-006]} Модуль «Освіта» повинен автоматично розраховувати інтегральні компетенції учнів на основі поєднання поточних балів за встановленою вагою.
    \item \textbf{[REQ-EDU-FUN-007]} Модуль «Освіта» повинен підтримувати створення та затвердження навчальних планів за допомогою регламентованого процесу погодження (Методист --- Завуч --- Директор).
    \item \textbf{[REQ-EDU-FUN-008]} Модуль «Освіта» повинен забезпечувати перевірку наявності колізій у розкладі занять (накладки викладачів, зайнятість аудиторій).
    \item \textbf{[REQ-EDU-FUN-009]} Модуль «Освіта» повинен надсилати сповіщення у разі фіксації критично низьких оцінок ($<4$) учням та їхнім батькам.
    \item \textbf{[REQ-EDU-FUN-010]} Модуль «Освіта» повинен забезпечувати вивантаження даних успішності до державної системи ЄДЕБО за допомогою захищеного HTTPS REST API.
\end{itemize}

\subsection{4.2. Вимоги до інтерфейсів та дизайну}
\begin{itemize}
    \item \textbf{[REQ-EDU-UI-001]} Користувацький інтерфейс повинен відповідати вимогам доступності ДСТУ EN 301 549:2022 (WCAG 2.1 AA) для забезпечення зручності роботи користувачів з особливими потребами.
    \item \textbf{[REQ-EDU-UI-002]} Веб-інтерфейс повинен бути реалізований за технологією Single Page Application (SPA) на базі WebSocket для мінімізації трафіку та забезпечення миттєвого відгуку.
    \item \textbf{[REQ-EDU-UI-003]} Система повинна використовувати підхід Data-on-Wire --- передавати тільки чисті структуровані дані у форматі JSON без повторного пересилання HTML-шаблонів.
    \item \textbf{[REQ-EDU-UI-004]} Дизайн системи повинен автоматично адаптуватися під екрани з розширенням від 320px до 3840px (Responsive Web Design).
\end{itemize}

\subsection{4.3. Вимоги до безпеки та захисту інформації}
\begin{itemize}
    \item \textbf{[REQ-EDU-SEC-001]} Автентифікація всіх посадових осіб (адміністратори, педагоги, методисти) повинна здійснюватися виключно за допомогою КЕП у форматі CAdES-X Long.
    \item \textbf{[REQ-EDU-SEC-002]} Передача даних у мережі повинна здійснюватися виключно через захищений канал HTTPS (TLS 1.3).
    \item \textbf{[REQ-EDU-SEC-003]} Усі відкриті порти, крім порту 443 (HTTPS), повинні бути закриті на рівні мережевого екрану.
    \item \textbf{[REQ-EDU-SEC-004]} Оскільки реляційні бази даних не використовуються, ризик SQL-ін'єкцій відсутній; натомість Erlang-сервіси повинні виконувати суворе очищення текстових полів для запобігання XSS-атакам.
    \item \textbf{[REQ-EDU-SEC-005]} Модуль авторизації повинен перевіряти ABAC-атрибути суб'єкта та об'єкта доступу перед виконанням будь-якої операції читання чи запису.
    \item \textbf{[REQ-EDU-SEC-006]} Токени користувацьких сесій повинні записуватися в Cookie з атрибутами \texttt{Secure}, \texttt{HttpOnly} та \texttt{SameSite=Strict}.
    \item \textbf{[REQ-EDU-SEC-007]} Модуль авторизації повинен блокувати паралельні сесії одного користувача з різних IP-адрес.
    \item \textbf{[REQ-EDU-SEC-008]} Дані користувачів повинні бути захищені відповідно до Закону «Про захист персональних даних».
\end{itemize}

\subsection{4.4. Вимоги до продуктивності та надійності}
\begin{itemize}
    \item \textbf{[REQ-EDU-PERF-001]} Час реакції системи на створення асинхронної фонової операції не повинен перевищувати 200 мс.
    \item \textbf{[REQ-EDU-PERF-002]} Пропускна здатність черги обробки транзакцій Mnesia повинна складати не менше 500 операцій на секунду на один серверний вузол.
    \item \textbf{[REQ-EDU-PERF-003]} Затримка оновлення даних в електронному журналі через Secure WebSockets не повинна перевищувати 150 мс.
    \item \textbf{[REQ-EDU-PERF-004]} Час повного рендерингу та відображення сторінки журналу на клієнті не повинен перевищувати 1 секунди при розмірі сторінки до 2000 записів.
\end{itemize}

\subsection{4.5. Вимоги до логування та аудиту}
\begin{itemize}
    \item \textbf{[REQ-EDU-LOG-001]} Усі критичні системні події (авторизація, зміна оцінок, накладання КЕП, помилки авторизації) повинні реєструватися у централізованому журналі аудиту (Audit Log).
    \item \textbf{[REQ-EDU-LOG-002]} Логи повинні експортуватися у режимі реального часу до стеку ELK (Elasticsearch, Logstash, Kibana) за допомогою Filebeat.
    \item \textbf{[REQ-EDU-LOG-003]} Запис логу повинен мати формат JSON та обов'язково містити поля: мітка часу (\texttt{time}), рівень критичності (\texttt{severity}), унікальний код події, ідентифікатор користувача (\texttt{user\_id}) та деталі помилки.
\end{itemize}

\subsection{4.6. Адміністративні та юридичні вимоги}
\begin{itemize}
    \item \textbf{[REQ-EDU-ADM-001]} Програмний код Модуля «Освіта» повинен бути вкритий автоматичними тестами (Unit та Integration) не менше ніж на 80\%.
    \item \textbf{[REQ-EDU-ADM-002]} Для тестування у непублічних середовищах повинна бути передбачена можливість запуску без підключення до реальних сервісів КЕП.
    \item \textbf{[REQ-EDU-ADM-003]} Вендор повинен надати повну технічну документацію на засіб інформатизації, що містить Настанову користувача, Настанову адміністратора та Опис API за стандартом ДСТУ 3008:2015.
    \item \textbf{[REQ-EDU-ADM-004]} Усі майнові права інтелектуальної власності на створене програмне забезпечення повинні бути повністю передані замовнику (державі в особі ДСА України).
\end{itemize}

\end{document}
